%=====================================================================
\chapter{Communicating Data Science}\label{ch:communication}
%=====================================================================
\locator{6}{Part VI --- Professional Practice}

\begin{outcomes}
By the end of this chapter you will be able to:
\begin{tight}
  \item \textbf{Adapt} the same result for a technical, managerial and affected-individual audience.
  \item \textbf{Structure} a findings report so the conclusion arrives first.
  \item \textbf{Express} uncertainty in language a non-specialist will act on correctly.
  \item \textbf{Design} a dashboard around a decision rather than around available data.
  \item \textbf{Recognise} the presentation choices that mislead, including your own.
\end{tight}
\end{outcomes}

\begin{prereq}
\chref{eda} (chart selection and anti-patterns), \chref{inference} (what an interval
does and does not mean) and \chref{ethics} (explaining a decision to the person it
affects).
\end{prereq}

\begin{keyterms}
audience analysis $\bullet$ bottom line up front $\bullet$ executive summary
$\bullet$ narrative arc $\bullet$ uncertainty communication $\bullet$ natural
frequencies $\bullet$ actionable insight $\bullet$ dashboard $\bullet$ decision
support $\bullet$ caveat
\end{keyterms}

\section{The Work Is Not Finished Until Someone Acts}

\begin{conceptbox}[The Uncomfortable Truth of This Chapter]
A correct analysis nobody understands has the same effect on the world as no
analysis. Communication is not a soft skill appended to the technical work --- it is
the step at which the technical work either produces value or does not. More good
data science dies in the final presentation than in modelling.
\end{conceptbox}

\section{Three Audiences, One Result}

\dsfig{%
\begin{tikzpicture}[font=\scriptsize,
  hd/.style={rounded corners=3pt, fill=#1, text=white, font=\tiny\bfseries,
             text width=4.15cm, align=center, minimum height=0.5cm},
  b/.style={rounded corners=3pt, draw=#1, fill=#1!7, text width=4.15cm, align=left,
            font=\tiny, inner sep=5pt, minimum height=3.4cm}]

\node[hd=secondary] at (0,0)    {TECHNICAL PEER};
\node[b=secondary]  at (0,-2.0) {%
  \textbf{Wants:} to check your work.\\[3pt]
  \textbf{Give:} method, validation scheme, PR-AUC with CV variance, ablations,
  known failure modes, code.\\[3pt]
  \textbf{Length:} as long as needed.\\[3pt]
  \emph{``Walk-forward CV, 5 folds; PR-AUC 0.41 $\pm$0.04 versus 0.22 baseline;
  recall 0.71 at a 20\% flag rate.''}};

\node[hd=goldhl] at (4.75,0)    {DECISION MAKER};
\node[b=goldhl]  at (4.75,-2.0) {%
  \textbf{Wants:} to decide something.\\[3pt]
  \textbf{Give:} what it does, what it costs, what it risks, what you recommend.
  One caveat that changes the decision.\\[3pt]
  \textbf{Length:} one page. Really.\\[3pt]
  \emph{``We can find 7 in 10 failures by week 4, contacting 1 in 5 students. Current
  practice finds fewer than half. Pilot needs 0.4 FTE.''}};

\node[hd=greenhl] at (9.5,0)    {AFFECTED INDIVIDUAL};
\node[b=greenhl]  at (9.5,-2.0) {%
  \textbf{Wants:} to know what to do.\\[3pt]
  \textbf{Give:} the specific reason, and an action that changes the outcome.
  Never a score, never a label.\\[3pt]
  \textbf{Length:} two sentences.\\[3pt]
  \emph{``You haven't submitted since week 2. Would a session with an adviser this
  week help?''}};

\node[anchor=west, align=left, text width=13.8cm, font=\tiny, text=primary!60!black]
     at (-2.2,-4.15)
     {\textbf{Same model, same evidence, three documents.} The commonest professional
      error is writing the left-hand column and sending it to all three audiences ---
      which reads as rigour to the author and as evasion to everyone else. The third
      column is the hardest and the most consequential: it is a person's only view of a
      system that is making a judgement about them (\chref{ethics}).};
\end{tikzpicture}}%
{One result, three audiences. What changes is not the honesty but the unit of
explanation: a metric for the peer, a decision for the manager, an action for the
person.}{three-audiences}

\section{Structure: Conclusion First}

\begin{definitionbox}[Bottom Line Up Front]
State the conclusion and the recommended action in the first two sentences, then
support it. Academic writing builds to a conclusion; professional writing leads with
it, because your reader may stop after the first paragraph and must still leave with
the right answer.
\end{definitionbox}

\begin{worked}[The Same Finding, Written Two Ways]
\textbf{Version A --- the way students write it.}

\emph{``We extracted VLE clickstream data for the 2025 cohort and performed an
audit, finding 32\% missingness in prior qualification. After cleaning and
engineering 40 features, we compared logistic regression, random forest and
gradient boosting using stratified 5-fold cross-validation. The gradient boosting
model achieved the highest PR-AUC of 0.41. We then tuned the decision threshold and
examined subgroup performance\dots''}

\smallskip
The reader is four sentences in and does not yet know whether anything works or what
they are being asked to do. A busy head of department stops here.

\smallskip
\textbf{Version B --- bottom line up front.}

\emph{``We can identify about 7 in 10 students who will fail, by week 4, while
contacting only 1 in 5 of the cohort. That is a substantial improvement on the
current rule, which finds fewer than half. We recommend a one-semester pilot with
three advisers, at an estimated cost of 0.4 FTE.}

\emph{Three caveats: the model is trained on one cohort and should not be assumed to
transfer; it identifies students who resemble past failures, not students who will
benefit most from help; and roughly 4 in 5 flagged students would have passed
anyway, so the outreach must be framed as an offer of support rather than a
warning.''}

\smallskip
\textbf{What changed.} The number that leads is the one tied to the decision. The
technical work is unchanged --- and is available in the appendix for the peer
audience. Note that Version B is also \emph{more} honest: the precision caveat is
stated plainly rather than buried, because a reader acting on this must know that
most flagged students are fine.
\end{worked}

\section{Communicating Uncertainty}

\begin{alertbox}[Percentages Mislead; Natural Frequencies Do Not]
``The model has 78\% precision'' is understood by almost nobody outside this
course. ``Of every 100 students we contact, about 78 would genuinely have failed''
is understood by everybody. The research on risk communication is consistent on
this: express probabilities as counts out of a stated population, and people reason
about them correctly.

Likewise, replace ``95\% CI [39.1, 44.9]'' with ``our best estimate is 42 minutes;
it would not surprise us if the true figure were anywhere between 39 and 45.''
\end{alertbox}

\rowcolors{2}{white}{lightbg}
\begin{longtable}{@{}p{5.6cm}p{8.4cm}@{}}
\toprule
\rowcolor{primary!15}
\textbf{Do not say} & \textbf{Say instead} \\
\midrule
\endhead
``The model is 94\% accurate.'' & ``It correctly classifies 94 in 100 --- but 97 in 100 cases are negative, so a model that always said `no' would score 97.'' \\
``$p < 0.05$, so it works.'' & ``The improvement is larger than we would expect from chance; the effect is about 11 percentage points.'' \\
``There is a 95\% chance the true value is in this range.'' & ``If we repeated this study many times, about 95 of every 100 such ranges would contain the true value.'' \\
``The model predicts this student will fail.'' & ``This student resembles past students of whom about 1 in 4 failed.'' \\
``The algorithm decided.'' & ``The model scored this case highly; the decision is yours, and here is why it scored highly.'' \\
\bottomrule
\end{longtable}

\begin{conceptbox}[Never Report a Point Estimate Alone]
A single number invites false confidence. Give a range, or give the number with the
one caveat that most changes how it should be used. This is not hedging --- a
stakeholder who later discovers an unmentioned uncertainty will discount everything
else you have told them.
\end{conceptbox}

\section{Dashboards}

\begin{alertbox}[Design for the Decision, Not for the Data]
The failure mode is a dashboard containing everything that was easy to compute. Ask
instead: \emph{who looks at this, how often, and what do they do differently as a
result?} If a chart cannot change an action, remove it. A dashboard with four
decision-relevant numbers is used; one with forty is glanced at once.
\end{alertbox}

\rowcolors{2}{white}{defbg}
\begin{longtable}{@{}p{3.0cm}p{4.5cm}p{6.5cm}@{}}
\toprule
\rowcolor{primary!15}
\textbf{Principle} & \textbf{Means} & \textbf{In practice} \\
\midrule
\endhead
One screen & No scrolling for the primary view & If it does not fit, you have more than one dashboard \\
Top-left first & Most important number in the reading-order corner & The single figure the viewer came for \\
Comparison always & A number alone means nothing & Versus target, versus last period, versus baseline \\
Consistent colour & One meaning per colour throughout & Red always means ``needs attention'', never merely ``category 2'' \\
Show uncertainty & Bands, not bare lines & A forecast without a range invites over-confidence \\
State the freshness & When was this last updated? & A stale dashboard is worse than none \\
\bottomrule
\end{longtable}

\section{Presentation Choices That Mislead}

\begin{pitfall}[Including the Ones You Will Be Tempted By]
\begin{tight}
  \item \textbf{Reporting the best of many runs.} If you tried twelve
        configurations, the best one is optimistically biased (\chref{inference}).
  \item \textbf{Quoting a metric without its threshold or baseline.} Both are
        required for the number to mean anything.
  \item \textbf{Choosing the metric after seeing the results.} Decide at framing
        time (\chref{lifecycle}).
  \item \textbf{Showing the successful segment.} ``It works well for full-time
        students'' when it was evaluated across all students is a selective report.
  \item \textbf{A truncated axis on a bar chart} (\chref{eda}) --- the most common
        way an honest person misleads accidentally.
  \item \textbf{Omitting the failure modes} because they weaken the story. They will
        be discovered in production instead, by someone with less context and less
        goodwill.
\end{tight}
\end{pitfall}

%---------------------------------------------------------------------
\begin{fourlenses}{Communicating Results}
\lensLA{Three audiences with genuinely different needs. The \emph{committee} needs
cohort-level effect and cost. The \emph{adviser} needs a ranked list with a
one-line reason per student, and the authority to disagree with it. The
\emph{student} needs a specific, actionable, non-stigmatising statement --- never a
risk score, and never the word ``predicted to fail''. Getting the third audience
right is an ethical requirement (\chref{ethics}), not a stylistic preference.}

\lensPM{Executives want a range and a confidence level, not a point estimate:
``we will complete 32--41 points, and we are more likely to land in the lower half
if the API dependency is not resolved this week.'' Note the reflexive effect ---
publishing a forecast changes the team's behaviour --- so state whether the number is
a prediction or a target. Confusing those two destroys the credibility of both.}

\lensIOT{The audience is a technician holding a tablet, not an analyst. The
deliverable is not a probability but an instruction: \emph{``Unit 47: bearing
vibration rising for 6 days. Inspect within 2 weeks. Evidence: 3$\times$ normal
variance since 14 July.''} Include the evidence, because a technician who can see
why will tell you when the model is wrong --- which is the feedback loop your
monitoring depends on (\chref{mlops}).}

\lensSEC{Two audiences, both hard. The \emph{analyst} needs context, evidence and a
recommended action within seconds --- alert fatigue is partly a communication design
problem, not only a threshold problem (\chref{cybersecurity}). The \emph{board}
needs the business case for work whose success is invisible; talk in terms of time
to detect, coverage and risk reduction, never model accuracy. Both need to know what
the system \emph{cannot} see, because an unstated blind spot becomes an assumed
capability.}
\end{fourlenses}

%---------------------------------------------------------------------
\begin{lab}[Rewrite a Result for Three Audiences]
Take one result from your own work --- or use the one below --- and write all three
versions. The exercise is harder than it looks and is worth doing before any real
presentation.

\begin{lstlisting}[basicstyle=\ttfamily\tiny, frame=none, backgroundcolor=\color{codebg}]
RAW RESULT
  Gradient boosting, walk-forward CV (5 folds).
  PR-AUC 0.41 (+/- 0.04) vs 0.22 majority baseline.
  At threshold 0.18: recall 0.71, precision 0.23, flag rate 0.20.
  Subgroups: recall 0.73 (full-time), 0.64 (part-time). Gap 9pp.

--- 1. TECHNICAL PEER (method + limits) ------------------------------
"Walk-forward 5-fold CV on 2024-25 cohorts; PR-AUC 0.41 +/- 0.04 against a
 0.22 majority baseline. Operating point chosen by cost ratio (FN:FP = 267:1),
 giving recall 0.71 at a 20% flag rate. Recall is 9pp lower for part-time
 students (n=88), which is within CV noise but is being monitored as a
 fairness risk. Features are strictly week-1-to-4; leakage review passed."

--- 2. DECISION MAKER (action + cost + caveat) -----------------------
"We can find about 7 in 10 students who will fail, by week 4, by contacting
 1 in 5 of the cohort. The current rule finds fewer than half. A pilot needs
 0.4 FTE of adviser time for one semester.
 One caveat you must plan around: about 4 of every 5 students we flag would
 have passed anyway, so this must be framed as an offer of support."

--- 3. AFFECTED INDIVIDUAL (specific, actionable, non-stigmatising) ---
"We noticed you haven't submitted since week 2, and students in that
 position often find the module hard to catch up on. Would a 20-minute
 session with an adviser this week be useful? Submitting this week's
 exercise would also help."
   -- no score, no "at risk" label, no "predicted to fail",
      and a concrete action the student can take.
\end{lstlisting}

\textbf{Check yourself:} version 3 contains no number at all, and is the version
that most affects an outcome.
\end{lab}

%---------------------------------------------------------------------
\begin{checkpoint}
\begin{tightnum}
  \item Rewrite ``the model achieves 82\% precision'' as a natural-frequency
        statement.
  \item Your dashboard has 22 charts. What single question decides which ones stay?
  \item Why should a report state the model's failure modes even when nobody asked?
\end{tightnum}
\end{checkpoint}

%---------------------------------------------------------------------
\begin{chaptersummary}
\begin{tight}
  \item An analysis nobody acts on has produced nothing; communication is the step
        that converts work into value.
  \item Technical peers need method and limits, decision makers need action and
        cost, affected individuals need a specific and actionable statement.
  \item Lead with the conclusion and the recommendation; support afterwards.
  \item Express probability as natural frequencies, and never report a point
        estimate without a range or the caveat that most changes its use.
  \item Design dashboards around a decision; remove anything that cannot change an
        action.
  \item State failure modes and blind spots proactively --- an unstated limitation
        becomes an assumed capability.
\end{tight}
\end{chaptersummary}

\begin{reviewq}
  \item \textbf{(MCQ)} ``Bottom line up front'' means: (a)~put the data first
        (b)~state the conclusion and recommendation first (c)~summarise at the end
        (d)~lead with the method.
  \item \textbf{(MCQ)} Which is a natural-frequency statement? (a)~``precision is
        0.78'' (b)~``$p<0.05$'' (c)~``about 78 of every 100 we contact would have
        failed'' (d)~``AUC 0.91''.
  \item \textbf{(Short)} Explain why a point estimate reported alone is a
        communication failure.
  \item \textbf{(Short)} Give two reasons never to show a student their own risk
        score.
  \item \textbf{(Short)} Why is alert fatigue partly a communication design problem?
  \item \textbf{(Applied)} Take the raw result in the lab above and write the three
        versions for a \emph{predictive maintenance} model instead, naming your three
        audiences.
  \item \textbf{(Applied)} A dashboard shows 18 charts and is never used. Redesign it:
        state who uses it, what decision they make, the four things you would keep,
        and what you would remove and why.
\end{reviewq}
